\chapter{Segurança e Conformidade}
\label{ch:seguranca}

Para a análise de segurança considerou-se que a plataforma seria comercializada como um brinquedo educativo, destinado a crianças a partir dos 8 anos, comandado pelo telemóvel. Esta hipótese é a mais exigente, porque os brinquedos têm requisitos próprios de segurança e o utilizador não tem formação técnica. Um protótipo académico, usado apenas em laboratório, não é colocado no mercado e não tem de ostentar a marcação \gls{ce} \cite{blueguide2022}; ainda assim, a análise permitiu incorporar no projeto medidas que um produto certificado teria de ter.

\section{Legislação e normas aplicáveis}
Como brinquedo elétrico com um emissor Bluetooth e uma bateria, o produto ficaria abrangido pelos atos e normas da tabela~\ref{tab:legislacao}.

\begin{table}[H]
\centering\small
\caption{Legislação e normas aplicáveis a um brinquedo com estas características.}
\begin{tabularx}{\textwidth}{>{\raggedright\arraybackslash}p{4.6cm}Y}
\toprule
Legislação ou norma & Âmbito\\\midrule
Diretiva 2009/48/CE \cite{toy2009} & Segurança dos brinquedos; em vigor\\
Regulamento (UE) 2025/2509 \cite{toy2025} & Substitui a diretiva a partir de 1 de agosto de 2030; introduz o passaporte digital do produto e exige avaliar os riscos dos brinquedos ligados digitalmente \cite{eurofinstoy}\\
Diretiva 2014/53/UE (\acrshort{red}) \cite{red2014} & Equipamentos de rádio (Bluetooth)\\
Diretivas 2011/65/UE (\acrshort{rohs}) e 2012/19/UE (\acrshort{reee}) \cite{rohs2011,weee2012} & Substâncias perigosas e fim de vida do equipamento elétrico\\
Regulamento (UE) 2023/1542 \cite{battery2023} & Baterias\\
EN 71-1, EN 71-2 e EN 71-3 \cite{en71} & Propriedades mecânicas e físicas, inflamabilidade e migração de elementos químicos\\
EN IEC 62115 \cite{en62115} & Segurança de brinquedos elétricos\\
EN 300 328, EN 301 489-1 e -17 \cite{en300328,en301489} & Espetro e compatibilidade eletromagnética do rádio\\
\bottomrule
\end{tabularx}
\label{tab:legislacao}
\end{table}

\section{Requisitos para brinquedos e resposta do projeto}
A tabela~\ref{tab:brinquedo} cruza os principais requisitos de segurança dos brinquedos com a situação do protótipo.

\begin{table}[H]
\centering\small
\caption{Requisitos de segurança para brinquedos e situação do protótipo.}
\begin{tabularx}{\textwidth}{>{\raggedright\arraybackslash}p{3.2cm}YY}
\toprule
Requisito & Situação no protótipo & Alteração para um produto\\\midrule
Tensão nominal não superior a \SI{24}{\volt} \cite{toy2009} & Cumpre: máximo de \SI{8.4}{\volt} (RNF12) & --\\
Bateria e eletrónica inacessíveis sem ferramenta & Cumpre: caixa fechada com parafusos (SEG-08) & --\\
Sem arestas vivas nem pontas (EN 71-1) & Peças impressas com arestas por acabar & Arredondar arestas e rever o acabamento\\
Sem pontos de entalamento de dedos (EN 71-1) & Folgas entre quadro, plataforma e torres; servos com binário elevado & Proteger as zonas de rotação, limitar velocidade e binário, manter o interruptor de corte dos servos\\
Temperatura das partes acessíveis (EN IEC 62115) & Servos e reguladores podem aquecer em bloqueio & Medir temperaturas; deteção de bloqueio em software\\
Bateria segura & LiPo sem proteção própria & Usar bateria com circuito de proteção e carregamento integrado, sem intervenção da criança\\
Materiais (EN 71-2, EN 71-3) & PLA sem tintas & Confirmar inflamabilidade e migração de elementos com o fornecedor do material\\
Brinquedo ligado & Recebe apenas ângulos; não recolhe dados pessoais & Ligação protegida por emparelhamento; informação ao utilizador\\
\bottomrule
\end{tabularx}
\label{tab:brinquedo}
\end{table}

O protótipo cumpre os requisitos elétricos de base, mas não está pronto para ser um brinquedo: as maiores lacunas são a bateria LiPo sem proteção e os pontos de entalamento no mecanismo. Ambas têm solução conhecida e baixo custo.

\section{Análise de riscos}
A análise seguiu o método da EN ISO 12100: identificar os perigos, estimar o risco como produto da gravidade (G) pela probabilidade (P), numa escala de 1 a 3, e definir medidas para o reduzir. Considera-se aceitável um risco residual igual ou inferior a 3.

\begin{table}[H]
\centering\scriptsize
\caption{Análise de riscos.}
\begin{tabularx}{\textwidth}{>{\raggedright\arraybackslash}p{2.4cm}ccc>{\raggedright\arraybackslash}Yc}
\toprule
Perigo & G & P & R & Medidas de redução & R res.\\\midrule
Entalamento de dedos & 2 & 2 & 4 & Limites de amplitude e velocidade; proteção das zonas de rotação; interruptor de corte dos servos & 2\\
Movimento inesperado & 2 & 2 & 4 & Arranque desarmado; armar explícito; \textit{watchdog}; modo seguro & 2\\
Incêndio da bateria & 3 & 1 & 3 & Fusível; bateria com proteção; monitorização da tensão & 2\\
Sobreaquecimento dos servos & 1 & 2 & 2 & Limites de comando; deteção de bloqueio & 1\\
Arestas vivas & 1 & 2 & 2 & Acabamento das peças & 1\\
Choque elétrico & 1 & 1 & 1 & Muito baixa tensão & 1\\
\bottomrule
\end{tabularx}
\label{tab:riscos}
\end{table}

\section{Funções de segurança implementadas}
As medidas anteriores traduzem-se nas seguintes funções, verificadas pelos testes T-SEG (capítulo~\ref{ch:testes}):
\begin{itemize}
\item \textbf{Corte dos servos por hardware:} um interruptor de dois polos corta a alimentação dos servos sem depender do microcontrolador.
\item \textbf{Arranque seguro:} o sistema arranca desarmado e só se move depois de calibrado e armado.
\item \textbf{Modo seguro:} sem referências válidas durante \SI{1}{\second}, ou com bateria baixa, a plataforma volta suavemente à horizontal.
\item \textbf{Limites em software:} amplitude, velocidade e largura de impulso dos servos limitadas.
\item \textbf{\textit{Watchdog}:} reinício automático em caso de bloqueio do programa.
\item \textbf{Proteção elétrica:} fusível, interruptor geral e eletrónica fechada numa caixa aparafusada.
\end{itemize}
